\documentclass{article}
\usepackage{amsmath}
\usepackage[hidelinks]{hyperref}
\title{Give the forecast an actual pre-mission consequence}
\date{}
\begin{document}
\maketitle

\paragraph{Two episodes and one fixed law.}
The paired-versus-step-result comparison concerns acquired
calibration data for a \emph{later} loaded job, not a diagnosis
supplied for free when the future job is already under way.
Materialize the actual archived early native-result frequencies
along with the pinned RWARE and SubT code in \verb|SOURCE.tsv|.
A single unknown law fixes the independently realized loaded
pallet rate at $1/4$ or $3/4$ in both the earlier calibration and
future jobs. A current initial lane $X$ is known from the job
start, not the future blockage. Prepare for staying or switching
lane, or wait centrally. The future latent obstacle is inserted
only after the chosen preparation has physically completed.
The planner checks actual native forward motion and takes the
previously authorized switch route if blocked. It does not
choose a convenient new forecast law after provoking an outcome.

\paragraph{Commitment and native consequences.}
One reserve carrier has the named rail of slots at column four,
not permission to enter the loaded carrier's corridor on column
three. It starts at the center bay facing right; preparing either
route takes an actual turn and move, and leaves it on that end
of the rail facing out. The prep claim holds its pose and heading
until the named landing cut. Voluntary early relocation or turning
is not in this reference authority. These are task constraints,
not inherent limits of the stock RWARE package. Neither policy
gets a second helper or uncharged access to a better staging
position. Standard stock actions implement all executed motions
and the separate side action.

At the scheduled landing cut the native loaded carrier's actual
position determines where the reserve's forwarding endpoint can
receive the old obligation. Exactly adjacent bays support the
modeled path. If the forecast was correct the real three-byte
callback can occur at the initial due cut. A wrong lane requires
two turns and two native moves after release before the callback;
the central wait requires one turn and one move. None waives
the old report. Preparation displaces an optional four-byte
callback usable only from the central bay at its earlier appointed
cut. Its failed attempt is charged in the native radio ledger
rather than ignored. Waiting at the central bay preserves it.
Once the slot is lost, later relocation cannot retroactively
create that callback. A separate independent physical action
and radio job run in all cases. The fixed activation seam, time
mapping and named claims are supplied by this joint simulator;
there is no measured radio range or physical experiment.

\paragraph{Use only the calibration result, then compare complete rules.}
The direct replay enumerates the three decisions for each actual
allowed calibration outcome and reads the observed prior archive
without access to the future latent variable. With no useful
calibration, the two laws imply equal-prior future detour
probability $1/2$; random prep has mean delay two just like wait,
and wait keeps the optional report. A nonerased informative
result gives equal-prior conditional preferred-route probability
$5/8$: there is a $3/4$ posterior chance of the matching law and
future route is still probabilistic under each law. Delay if
staging the preferred route therefore averages $4(3/8)=3/2$.
Waiting costs two actual motion steps at that cut. If a lost
optional callback has normalized cost $c$ motion-step units,
the threshold is exactly $c<1/2$ to prepare after an informative
sample. At equality waiting is a valid tie. After an erased
sample there is no forecast advantage and the optimal rule can
wait. A policy that stages after every sample pays the optional
loss on erased trials without improving expected correction.
The physical and radio costs are separate observed quantities;
$c$ is a declared preference, not measured from devices.

With independent half erasure the rule prepares only nonerased
cases if $c<1/2$; for example at $c=1/4$ worst-law cost is $15/8$,
versus $7/4$ with no erasure and $2$ from uniform waiting.
Equal-prior expected utility lower-bounds all worst-law policies,
including randomized ones, and the chosen mirrored rules have
equal realized expectation under the two laws. The exact
direct enumeration agrees. The value is conditional on the
scope of the archived action information: a differently
authorized future fallback would change the route law, so the
same measured bit would not select a correct preparation.
Neither the delay nor the rate estimates transfer to such a change.
This is a finite physically executed consequence for R3b, not
an acceptance claim or a generic improvement over
\href{https://arxiv.org/abs/1205.0858}{controlled sensing}.
\end{document}
